% ============================================================
% 第 18 章 同步发电机的并联运行（教材第五篇，印刷页 204--214）
% 节名：第一节 同步发电机的并网条件和方法；第二节 同步发电机电磁功率与功角特性；
%       第三节 同步发电机并联运行时的功率调节
% ============================================================
\chapter{同步发电机的并联运行}

电力系统中的同步发电机都是并联运行的。并网要解决三个问题：\textbf{怎么安全地并上去}（并网条件与方法）、\textbf{并上以后能承担多大负载}（功角特性与静态稳定）、\textbf{怎么分配有功与无功}（功率调节）。这一章是同步电机理论中最"工程"的一章，也是电力系统课程的接口。

\section{同步发电机的并网条件和方法的}

\subsection{理想并网条件}
把同步发电机投入电网并联运行，必须满足下列条件（否则并网瞬间会产生巨大的冲击电流与转矩，甚至损坏机组）：
\begin{enumerate}
  \item 发电机的\textbf{相序}与电网相序相同；
  \item 发电机的\textbf{频率}与电网频率相同；
  \item 发电机的\textbf{端电压大小}与电网电压相同；
  \item 发电机端电压的\textbf{相位}与电网电压相位一致（即电压相量重合）。
\end{enumerate}
条件 1—4 全部满足时，发电机与电网之间没有环流，合闸瞬间冲击最小，称为\textbf{理想并网}。

\begin{center}
\begin{tabular}{@{}p{2.2cm}p{4.4cm}p{5.2cm}@{}}
\toprule
条件 & 不满足时的后果 & 调整手段 \\
\midrule
相序不同 & 相当于两相短路，冲击电流极大 & 并网前核对接线（唯一必须先满足的） \\
频率不同 & 产生拍频，功率振荡、电流冲击 & 调原动机转速 \\
电压大小不等 & 合闸瞬间产生无功冲击电流 & 调励磁电流 \\
相位不一致 & 产生很大的有功冲击与转矩冲击 & 调相位（准同期最关键的"合闸时机"） \\
\bottomrule
\end{tabular}
\end{center}

\subsection{准同期法（准同步法）}
发电机先由原动机拖到接近同步转速，加励磁使端电压接近电网电压，然后\textbf{在相位一致的瞬间合闸}。
\begin{itemize}
  \item \textbf{灯光熄灭法}：把三组灯接在发电机与电网的同名相之间，灯全暗时表示两侧电压同相位，此时合闸。
  \item \textbf{灯光旋转法}：把一组灯接在一侧的同名相与另一侧的异名相之间，灯依次明暗旋转，表示频率有微小差；当直接接同相的那组灯熄灭、另两组灯亮度相同时合闸。
  \item 现代机组普遍采用\textbf{自动准同期装置}，自动调节频率与电压并捕捉合闸时机。
\end{itemize}
准同期法的优点是合闸冲击小，缺点是操作要求高、并网较慢。

\subsection{自同期法}
发电机不加励磁、转速接近同步转速时先合闸，随后立即投入励磁，靠定子磁场把转子\textbf{牵入同步}。优点：操作简单、并网迅速，可在系统故障时快速投入；缺点：合闸瞬间从电网吸取很大的励磁涌流（可达额定电流数倍）、引起电网电压短时下降。适用于事故备用机组的紧急并网。

\corepoint{并网条件的"代价排序"要记住：\textbf{相序错}是灾难（相当于短路），必须事先核对；\textbf{相位差}直接决定冲击大小（$\Delta U=2U\sin(\Delta\delta/2)$，差 $30^\circ$ 就有 $0.52U$ 的电压差），所以准同期法把"捕捉相位"当成核心；频率差与电压差相对好调整。考试若问"并网时为什么先调频后调压"，答：调频决定相位漂移速度，必须先把频差压到很小，才可能捕捉到相位一致的瞬间。}

\section{同步发电机电磁功率与功角特性}

第 17 章已得隐极机的功角特性
\begin{equation}
P_{em}=\frac{mE_0U}{X_t}\sin\delta .
\end{equation}
这一节讨论它的工程含义。

\subsection{功角的物理意义}
$\delta$ 是励磁电动势 $\dot E_0$ 与端电压 $\dot U$ 之间的相位差，对应着\textbf{转子磁极轴线}与\textbf{气隙合成磁场轴线}之间的空间夹角。发电机输出有功功率越大，转子磁极必须越"超前"于气隙磁场——就像两人拉一根弹性绳，拉的力越大，绳子的伸长（角度差）越大。所以 $\delta$ 本质上是"电磁弹簧"的变形量。

\subsection{静态稳定与整步功率}
当 $\delta<90^\circ$ 时，$\dfrac{\mathrm{d}P_{em}}{\mathrm{d}\delta}>0$：负载增加时电机能自动增加电磁功率以平衡，运行是\textbf{稳定}的。
当 $\delta>90^\circ$ 时，$\dfrac{\mathrm{d}P_{em}}{\mathrm{d}\delta}<0$：负载再增加，电磁功率反而减小，转子继续加速、$\delta$ 继续增大——电机\textbf{失步}（失去同步）。

因此 $\delta=90^\circ$ 是隐极机的\textbf{静态稳定极限}，此时
\begin{equation}
P_{\max}=\frac{mE_0U}{X_t}
\end{equation}
称为\textbf{功率极限}。整步功率系数（比整步功率）定义为
\begin{equation}
P_{sy}=\frac{\mathrm{d}P_{em}}{\mathrm{d}\delta}=\frac{mE_0U}{X_t}\cos\delta ,
\end{equation}
它表征电机维持同步的能力：$\delta$ 越大 $P_{sy}$ 越小，越接近失步。

\keyformula{P_{\max}=\frac{mE_0U}{X_t},\qquad P_{sy}=\frac{mE_0U}{X_t}\cos\delta}

\subsection{过载能力}
定义过载能力（静态过载倍数）
\begin{equation}
k_m=\frac{P_{\max}}{P_N}=\frac{1}{X_t^{*}\sin\delta_N},
\end{equation}
一般要求 $k_m=1.7\sim2.0$，即额定运行时的功角 $\delta_N\approx25^\circ\sim35^\circ$，留出足够的稳定裕量。

\subsection{提高静态稳定的措施}
由 $P_{\max}=mE_0U/X_t$ 可见：
\begin{itemize}
  \item \textbf{增大励磁}（提高 $E_0$）可提高功率极限——这是励磁调节（自动电压调节器 AVR）稳定电网的基本手段；
  \item \textbf{减小同步电抗}（设计上增大气隙、增大短路比）可提高功率极限，但造价上升；
  \item 采用\textbf{快速励磁}、\textbf{电力系统稳定器（PSS）}可进一步改善动态稳定。
\end{itemize}

\section{同步发电机并联运行时的功率调节}

并网后的发电机与电网"锁"在一起，转速由电网频率决定，不能随意改变。此时调节的是\textbf{有功}与\textbf{无功}两个通道，且两者互相独立（这是同步电机最优雅的性质之一）。

\subsection{有功功率的调节——调原动机}
增大原动机（汽轮机/水轮机）的进汽量或进水量，输入机械功率增加，转子加速、$\delta$ 增大，电磁功率随之增大，直到与输入功率重新平衡。故
\begin{equation}
P_{em}\uparrow\ \Longleftrightarrow\ \delta\uparrow\ \Longleftrightarrow\ \text{原动机输入转矩}\uparrow .
\end{equation}
\textbf{有功功率只由原动机的输入功率决定，与励磁电流无关}。这就是电力系统"调频靠调有功"（调机组出力）的依据。

\subsection{无功功率的调节——调励磁}
在保持原动机输入不变（有功不变）的条件下增大励磁电流 $I_f$：$E_0$ 增大，功率因数角改变，电枢电流的无功分量增大，发电机向电网输出\textbf{感性无功}（过励状态）；反之减小励磁，发电机吸收感性无功（欠励状态）。

\textbf{无功功率只由励磁电流决定，与输入机械功率基本无关}。这是电力系统"调压靠调无功"（调励磁）的依据。

\begin{center}
\begin{tabular}{@{}p{2.4cm}p{3.4cm}p{3.8cm}p{3.2cm}@{}}
\toprule
调节量 & 调节对象 & 主要变化 & 电网侧效果 \\
\midrule
有功功率 & 原动机进汽/进水量 & $\delta$ 变化，$E_0$ 基本不变 & 调频（有功平衡） \\
无功功率 & 励磁电流 $I_f$ & $E_0$ 变化，$\delta$ 基本不变 & 调压（无功平衡） \\
\bottomrule
\end{tabular}
\end{center}

\subsection{并联运行时的无功分配}
多台发电机并联时，无功功率按各机的\textbf{励磁状态}分配，而非按容量自动分配：励磁电流偏大的机组承担更多无功。若励磁分配不当，会出现"一台过励、一台欠励"的环流，故需配置无功分配装置（无功电流补偿）。

\subsection{调相运行}
若把发电机轴上负载去掉、只带励磁运行，则电机成为\textbf{调相机}：过励时向电网输出感性无功（相当于并联电容器），欠励时吸收感性无功（相当于并联电抗器）。调相机用于长距离输电线路的无功补偿与电压支撑。

\corepoint{本章的两条"独立性"是全章的灵魂：\textbf{有功只跟原动机走，无功只跟励磁走}。由此可推出所有定性结论：增励磁 $\Rightarrow$ 无功增大、功率因数下降（过励）、端电压升高、功角略减；增原动机输入 $\Rightarrow$ 有功增大、$\delta$ 增大、无功基本不变。凡是"调某个量会怎样"的问题，先判断它属于有功通道还是无功通道，再作答，几乎不会错。}
